\documentclass[12pt,a4paper,openany]{book}
\usepackage[a4paper,top=2.5cm,bottom=2.5cm,left=3cm,right=2.5cm]{geometry}
\usepackage{fontspec}
\usepackage[brazil]{babel}
\usepackage{amsmath,amssymb}
\usepackage{microtype}
\usepackage{xcolor}
\usepackage{hyperref}
\usepackage{fancyhdr}
\usepackage{titlesec}
\usepackage{tcolorbox}
\tcbuselibrary{breakable,skins}
\setmainfont{Latin Modern Roman}
\setsansfont{Latin Modern Sans}
\definecolor{azul}{HTML}{0B4F7A}
\definecolor{laranja}{HTML}{EF6C00}
\hypersetup{colorlinks=true,linkcolor=azul,urlcolor=azul,pdftitle={Matemática Discreta - Livro de Revisão},pdfauthor={Jeann Victor}}
\titleformat{\chapter}[display]{\normalfont\huge\bfseries\color{azul}}{\chaptertitlename\ \thechapter}{12pt}{\LARGE}
\titleformat{\section}{\Large\bfseries\color{azul}}{\thesection}{0.8em}{}
\titleformat{\subsection}{\large\bfseries\color{laranja}}{\thesubsection}{0.8em}{}
\newtcolorbox{exemplo}[1]{enhanced,breakable,colback=gray!7,colframe=azul,title=#1,coltitle=white,fonttitle=\bfseries,arc=1mm,boxrule=.5pt}
\usepackage{enumitem}
\setlist{itemsep=2pt,topsep=4pt}
\setlength{\parindent}{1.25cm}
\setlength{\parskip}{0.15em}
\emergencystretch=4em
\sloppy
\raggedbottom
\pagestyle{fancy}
\fancyhf{}
\fancyhead[LE]{\small\thepage}
\fancyhead[RO]{\small\thepage}
\fancyhead[LO,RE]{\small Matemática Discreta}
\renewcommand{\headrulewidth}{0.4pt}
\setlength{\headheight}{15pt}
\begin{document}
\begin{titlepage}
\centering
\vspace*{3cm}
{\Huge\bfseries\color{azul} Matemática Discreta\par}
\vspace{1cm}
{\LARGE\color{laranja} Lógica, provas, conjuntos, indução, relações, funções e combinatória\par}
\vspace{2cm}\rule{0.65\textwidth}{0.4pt}\par\vspace{1.5cm}
{\Large Notas de aula organizadas a partir da disciplina de Matemática Discreta\par}
\vspace{1cm}
{\large Compiladas e organizadas por\par}\vspace{0.3cm}
{\Large\bfseries Jeann Victor\par}
\vfill
{\large Curso de Bacharelado em Ciência da Computação\par}
{\large 2025\par}
\end{titlepage}
\frontmatter
\chapter*{Prefácio}\addcontentsline{toc}{chapter}{Prefácio}
Este livro consolida os tópicos das dez listas da disciplina e transforma o material disperso em um roteiro de estudo coerente. O objetivo não é reproduzir os enunciados, mas explicar as ideias, destacar erros frequentes e mostrar estratégias de solução. As digitalizações originais e suas transcrições continuam disponíveis para conferência.
\tableofcontents
\mainmatter
\chapter{Fundamentos de lógica proposicional}



\section{Proposições e valor lógico}



Uma proposição é uma sentença declarativa à qual se pode atribuir exatamente um valor lógico: verdadeiro ou falso. Perguntas, ordens e expressões abertas não são proposições. A frase “7 é primo” é uma proposição verdadeira; “feche a porta” não possui valor lógico; e “x é par” só se torna proposição depois que o valor de x é fixado ou quantificado.



Essa distinção parece simples, mas evita um erro recorrente: tentar aplicar conectivos lógicos a frases cujo significado ainda depende de contexto. Ao formalizar um argumento, o primeiro passo é definir proposições atômicas claras. Por exemplo, p pode representar “o programa termina” e q pode representar “a entrada é válida”.



\section{Conectivos}



Os principais conectivos são negação, conjunção, disjunção, implicação e bicondicional. A negação de p, escrita $\neg p$, inverte seu valor lógico. A conjunção $p\land q$ é verdadeira somente quando ambas as parcelas são verdadeiras. A disjunção inclusiva $p\lor q$ é falsa somente quando ambas são falsas.



A implicação $p\to q$ é falsa apenas quando p é verdadeira e q é falsa. Isso reflete o compromisso expresso pela frase “se p, então q”: a única violação ocorre quando a condição se realiza, mas a consequência prometida não. A bicondicional $p\leftrightarrow q$ é verdadeira quando p e q têm o mesmo valor lógico.



\begin{exemplo}{Exemplo.}
Se p significa “n é múltiplo de 4” e q significa “n é par”, então $p\to q$ é verdadeira para todo inteiro n. A recíproca $q\to p$ é falsa: 6 é par, mas não é múltiplo de 4.
\end{exemplo}



\section{Equivalências úteis}



Duas fórmulas são logicamente equivalentes quando possuem a mesma coluna final em todas as linhas da tabela-verdade. Algumas equivalências aparecem continuamente em provas e simplificações:



\begin{itemize}
\item dupla negação: $\neg(\neg p)\equiv p$;
\item implicação: $p\to q\equiv\neg p\lor q$;
\item contrapositiva: $p\to q\equiv\neg q\to\neg p$;
\item leis de De Morgan: $\neg(p\land q)\equiv\neg p\lor\neg q$ e $\neg(p\lor q)\equiv\neg p\land\neg q$;
\item bicondicional: $p\leftrightarrow q\equiv(p\to q)\land(q\to p)$.
\end{itemize}



As leis de De Morgan mostram que, ao negar uma expressão composta, o conectivo muda e cada componente é negado. Em programação, isso ajuda a revisar condições complexas e a evitar inversões incorretas.



\section{Quantificadores e negação}



O quantificador universal $\forall$ afirma que uma propriedade vale para todos os elementos do domínio. O quantificador existencial $\exists$ afirma que existe pelo menos um elemento para o qual a propriedade vale. A ordem dos quantificadores importa: $\forall x\,\exists y\,P(x,y)$ não significa o mesmo que $\exists y\,\forall x\,P(x,y)$.



Negar um quantificador troca seu tipo e nega a propriedade:



\[\neg(\forall x\,P(x))\equiv\exists x\,\neg P(x),\]



\[\neg(\exists x\,P(x))\equiv\forall x\,\neg P(x).\]



\begin{exemplo}{Exemplo.}
A negação de “todo arquivo possui backup” é “existe pelo menos um arquivo que não possui backup”. Não é correto dizer “nenhum arquivo possui backup”.
\end{exemplo}



\section{Exercícios de fixação}



\begin{itemize}
\item Classifique como proposição ou não: “3 é ímpar”, “qual é sua idade?” e “x+1=5”.
\item Escreva a contrapositiva de “se um grafo é uma árvore, então ele é conexo”.
\item Negue: “para todo usuário existe uma senha válida”.
\item Verifique por tabela-verdade que $p\to q$ equivale a $\neg q\to\neg p$.
\end{itemize}



\chapter{Linguagem matemática e técnicas de demonstração}



\section{Axiomas, conjecturas e teoremas}



Um axioma é aceito como ponto de partida dentro de uma teoria. Uma conjectura é uma afirmação apoiada por evidências, mas ainda sem demonstração. Um teorema é uma afirmação demonstrada a partir de definições, axiomas e resultados anteriores. Um lema é um resultado auxiliar usado na prova de outro resultado; um corolário decorre de modo quase imediato de um teorema.



Uma demonstração não é apenas uma sequência de contas. Ela deve declarar hipóteses, indicar o objetivo e justificar cada passagem. Exemplos numéricos podem sugerir um resultado, mas não provam uma afirmação universal.



\section{Prova direta}



Para provar uma implicação $p\to q$ diretamente, assume-se p e, usando definições e resultados conhecidos, deduz-se q.



\begin{exemplo}{Exemplo.}
Provar que a soma de dois inteiros pares é par. Se $m=2a$ e $n=2b$, com a e b inteiros, então $m+n=2a+2b=2(a+b)$. Como $a+b$ é inteiro, $m+n$ é par.
\end{exemplo}



Um cuidado importante é não reutilizar a mesma variável para representar testemunhas independentes. Escrever $m=2k$ e $n=2k$ força indevidamente $m=n$. O correto é usar, por exemplo, $m=2a$ e $n=2b$.



\section{Contrapositiva e contradição}



Como $p\to q$ equivale a $\neg q\to\neg p$, às vezes é mais fácil provar a contrapositiva. Para mostrar “se $n^2$ é par, então n é par”, pode-se provar que, se n é ímpar, então $n^2$ também é ímpar.



Na prova por contradição, assume-se a negação do objetivo e mostra-se que essa hipótese leva a uma impossibilidade. A contradição deve surgir de modo explícito: uma afirmação e sua negação, uma violação de hipótese ou um resultado incompatível com uma propriedade conhecida.



\section{Como revisar uma prova}



\begin{itemize}
\item As variáveis e seus domínios foram declarados?
\item Cada hipótese foi usada de forma legítima?
\item O argumento prova a afirmação, ou apenas a repete com outras palavras?
\item Um caso particular foi confundido com uma prova geral?
\item A conclusão obtida é exatamente a conclusão pedida?
\end{itemize}



\begin{exemplo}{Erro clássico.}
Começar uma prova de “se p, então q” assumindo q. Isso pode produzir raciocínio circular. Em uma prova direta, assume-se p; q deve ser deduzida.
\end{exemplo}



\chapter{Conjuntos e operações}



\section{Elementos, subconjuntos e conjunto das partes}



Escrevemos $x\in A$ quando x é elemento de A e $A\subseteq B$ quando todo elemento de A também pertence a B. As afirmações $x\in A$ e $\{x\}\subseteq A$ são relacionadas, mas não idênticas: a primeira compara elemento e conjunto; a segunda compara dois conjuntos.



O conjunto vazio $\varnothing$ não possui elementos. Já $\{\varnothing\}$ possui um elemento: o próprio conjunto vazio. Para um conjunto finito com n elementos, seu conjunto das partes $\mathcal{P}(A)$ possui $2^n$ subconjuntos.



\section{União, interseção e diferença}



A união $A\cup B$ reúne os elementos que pertencem a A ou a B. A interseção $A\cap B$ contém os elementos comuns. A diferença $A\setminus B$ contém os elementos de A que não pertencem a B. Dado um universo U, o complemento de A é $A^c=U\setminus A$.



As leis de De Morgan para conjuntos são:



\[(A\cup B)^c=A^c\cap B^c, \qquad (A\cap B)^c=A^c\cup B^c.\]



Elas podem ser provadas pelo método de dupla inclusão: demonstra-se primeiro que o conjunto da esquerda está contido no da direita e depois a inclusão contrária.



\section{Produto cartesiano}



O produto cartesiano $A\times B$ é o conjunto de pares ordenados $(a,b)$ com $a\in A$ e $b\in B$. A ordem importa: em geral, $A\times B\neq B\times A$. Se A e B são finitos, então $|A\times B|=|A|\,|B|$.



\begin{exemplo}{Exemplo.}
Se $A=\{1,2\}$ e $B=\{x,y,z\}$, então $A\times B$ possui seis pares. O par $(1,x)$ não é igual a $(x,1)$.
\end{exemplo}



\section{Estratégia para identidades de conjuntos}



Para provar $A\times(B\cap C)=(A\times B)\cap(A\times C)$, escolha um par arbitrário $(x,y)$ no lado esquerdo. Então $x\in A$ e $y\in B\cap C$, logo y pertence a B e a C. Assim, $(x,y)$ pertence simultaneamente a $A\times B$ e a $A\times C$. A inclusão oposta segue revertendo os passos.



\chapter{Indução matemática}



\section{Estrutura da prova}



O princípio da indução finita é usado para provar uma propriedade $P(n)$ para todos os inteiros a partir de um valor inicial. A prova tem duas partes inseparáveis:



\begin{itemize}
\item base: verificar $P(n_0)$;
\item passo indutivo: assumir $P(k)$ para um k arbitrário e provar $P(k+1)$.
\end{itemize}



A hipótese de indução não afirma que a propriedade é verdadeira para todos os números; ela permite usar $P(k)$ temporariamente para construir o próximo caso.



\section{Soma dos números ímpares}



Vamos provar que $1+3+\cdots+(2n-1)=n^2$. Para $n=1$, ambos os lados valem 1. Suponha agora que a soma dos k primeiros ímpares seja $k^2$. Ao acrescentar o próximo ímpar, obtemos



\[k^2+[2(k+1)-1]=k^2+2k+1=(k+1)^2.\]



Portanto, a fórmula vale para todo $n\geq1$.



\section{Divisibilidade}



Para provar que $n^3-n$ é divisível por 3, há uma demonstração algébrica curta:



\[n^3-n=n(n-1)(n+1).\]



Esse é o produto de três inteiros consecutivos, e um deles é múltiplo de 3. A observação também ajuda a construir uma prova por indução: a diferença entre os casos consecutivos é



\[[(k+1)^3-(k+1)]-(k^3-k)=3k(k+1),\]



que é múltipla de 3.



\section{Sequências recorrentes}



Se $a_1=3$ e $a_k=7a_{k-1}$ para $k\geq2$, a forma fechada é $a_n=3\cdot7^{n-1}$. A base é imediata. No passo, se $a_k=3\cdot7^{k-1}$, então



\[a_{k+1}=7a_k=7\cdot3\cdot7^{k-1}=3\cdot7^k.\]



\section{Diagnóstico de uma indução incompleta}



\begin{itemize}
\item Sem a base, a cadeia lógica pode começar em lugar nenhum.
\item Sem usar a hipótese de indução, o argumento normalmente é uma prova direta disfarçada ou possui uma lacuna.
\item O passo deve provar o caso $k+1$, não repetir o caso k.
\item Se a recorrência depende de dois termos anteriores, pode ser necessário assumir dois casos ou usar indução forte.
\end{itemize}



\chapter{Relações binárias}



\section{Definição e representação}



Uma relação de A em B é qualquer subconjunto de $A\times B$. Quando A=B, dizemos que a relação é “em A”. Relações finitas podem ser representadas por pares ordenados, diagramas de setas, matrizes ou grafos dirigidos.



Uma relação R em A é reflexiva quando $xRx$ para todo x; simétrica quando $xRy$ implica $yRx$; antissimétrica quando $xRy$ e $yRx$ implicam x=y; e transitiva quando $xRy$ e $yRz$ implicam $xRz$.



\section{Relações de equivalência}



Uma relação de equivalência é reflexiva, simétrica e transitiva. Ela particiona o conjunto em classes de equivalência. Congruência módulo m é o exemplo central:



\[a\equiv b\pmod m \quad\Longleftrightarrow\quad m\mid(a-b).\]



Os inteiros que deixam o mesmo resto na divisão por m pertencem à mesma classe.



\section{Ordens parciais}



Uma ordem parcial é reflexiva, antissimétrica e transitiva. Nem todos os pares precisam ser comparáveis. A relação de inclusão $\subseteq$ no conjunto das partes é uma ordem parcial. Já $\leq$ nos reais é uma ordem total, pois quaisquer dois reais podem ser comparados.



\begin{exemplo}{Armadilha.}
Simetria e antissimetria não são opostos. Uma relação pode ser simultaneamente simétrica e antissimétrica, como a igualdade.
\end{exemplo}



\section{Fechos}



O fecho reflexivo adiciona todos os pares $(x,x)$ ausentes. O fecho simétrico adiciona $(y,x)$ sempre que $(x,y)$ aparece. O fecho transitivo adiciona as conexões exigidas por caminhos de comprimento maior que um. Em computação, o fecho transitivo modela alcançabilidade.



\chapter{Funções}



\section{Definição}



Uma função $f:A\to B$ associa a cada elemento do domínio A exatamente um elemento do contradomínio B. O conjunto imagem é formado pelos valores realmente atingidos. Domínio, contradomínio e regra de associação fazem parte da definição; duas funções com a mesma fórmula podem ser diferentes se seus domínios ou contradomínios diferem.



\section{Injetividade, sobrejetividade e bijetividade}



Uma função é injetiva quando entradas distintas produzem saídas distintas. Equivalentemente, $f(x_1)=f(x_2)$ implica $x_1=x_2$. Ela é sobrejetiva quando todo elemento do contradomínio é imagem de algum elemento do domínio. É bijetiva quando possui as duas propriedades.



\begin{exemplo}{Exemplo.}
A função $f:\mathbb{R}\to\mathbb{R}$, $f(x)=x+2$, é bijetiva. A igualdade $f(x_1)=f(x_2)$ implica $x_1=x_2$, provando injetividade. Para qualquer $y\in\mathbb{R}$, escolher $x=y-2$ produz $f(x)=y$, provando sobrejetividade.
\end{exemplo}



\section{Composição e inversa}



Se $f:A\to B$ e $g:B\to C$, então $(g\circ f)(x)=g(f(x))$. A ordem é essencial: em geral, $g\circ f\neq f\circ g$. Uma função admite inversa funcional se e somente se é bijetiva.



Para encontrar a inversa de $f(x)=ax+b$, com $a\neq0$, escreva $y=ax+b$ e isole x:



\[x=\frac{y-b}{a}.\]



Logo, $f^{-1}(y)=(y-b)/a$, respeitando os domínios apropriados.



\section{Funções definidas por partes}



Ao compor funções definidas por partes, a condição deve ser aplicada ao argumento interno. Se $f(u)$ muda de regra conforme $u\leq3$ ou $u>3$, então em $f(g(x))$ é necessário resolver quando $g(x)\leq3$ e quando $g(x)>3$.



\chapter{Princípios de contagem e combinatória}



\section{Princípios aditivo e multiplicativo}



Se uma escolha pode ser feita de a maneiras ou, em caso disjunto, de b maneiras, existem $a+b$ possibilidades. Se um processo possui etapas sucessivas independentes com a e b opções, existem $ab$ resultados.



\begin{exemplo}{Exemplo.}
Um teste de 10 questões com cinco alternativas possui $5^{10}$ gabaritos possíveis: para cada uma das dez posições há cinco escolhas.
\end{exemplo}



\section{Permutações e arranjos}



O número de ordenações de n objetos distintos é $n!$. Quando há repetições com multiplicidades $n_1,\ldots,n_k$, o número de sequências distintas é



\[\frac{n!}{n_1!\cdots n_k!}.\]



Escolher e ordenar k elementos distintos dentre n produz o arranjo



\[A(n,k)=\frac{n!}{(n-k)!}.\]



\section{Combinações}



Quando a ordem não importa, usamos



\[\binom{n}{k}=\frac{n!}{k!(n-k)!}.\]



Uma forma segura de decidir entre arranjo e combinação é perguntar se trocar a ordem dos selecionados cria um resultado diferente. Escolher presidente, vice e secretário exige ordem; escolher três representantes para uma cerimônia não.



\section{Agrupamentos e divisões}



Para dividir 12 pessoas em dois grupos não rotulados de 6, escolher o primeiro grupo e depois dividir por 2 corrige a troca dos grupos:



\[\frac{1}{2}\binom{12}{6}.\]



Em grupos rotulados, como “equipe A” e “equipe B”, essa divisão por 2 não é feita.



\section{Inclusão-exclusão}



Para conjuntos finitos,



\[|A\cup B|=|A|+|B|-|A\cap B|.\]



O termo da interseção é subtraído porque foi contado duas vezes. Com três conjuntos, somam-se os tamanhos individuais, subtraem-se as interseções dois a dois e soma-se novamente a interseção tripla.



\chapter{Estratégias para resolução de listas}



\section{Traduzir antes de calcular}



Em lógica, defina as proposições atômicas antes de manipular símbolos. Em conjuntos, determine o universo e o significado de cada operação. Em funções, registre domínio e contradomínio. Em combinatória, identifique se a ordem importa, se há repetição e se os grupos são rotulados.



\section{Testar casos pequenos}



Casos pequenos não substituem uma prova, mas ajudam a detectar fórmulas incorretas. Para uma afirmação sobre subconjuntos, teste conjuntos com zero, um e dois elementos. Para recorrências, calcule os primeiros termos. Para identidades combinatórias, compare os dois lados para valores pequenos de n.



\section{Trabalhar de trás para frente}



Quando o objetivo é uma igualdade, examine a forma final desejada e identifique qual transformação intermediária a produziria. Em indução, escreva o caso $k+1$ e tente separar uma parcela correspondente ao caso k. Em provas de conjuntos, traduza o pertencimento ao conjunto final em condições lógicas.



\section{Verificação dimensional e de domínio}



Uma resposta combinatória deve ser inteira e não negativa. Uma inversa precisa respeitar domínio e imagem. Uma expressão como $\binom{n}{k}$ exige $0\leq k\leq n$. A checagem de domínio elimina muitos erros antes mesmo da revisão algébrica.



\chapter{Revisão integrada}



\section{Problema 1 - negação com quantificadores}



Negue a afirmação: “para todo inteiro n existe um inteiro m maior que n e primo”. A estrutura é $\forall n\,\exists m\,P(n,m)$. A negação correta é $\exists n\,\forall m\,\neg P(n,m)$: “existe um inteiro n tal que nenhum inteiro m maior que n é primo”. A afirmação negada é falsa, mas sua forma lógica está correta.



\section{Problema 2 - bijeção}



Considere $f:\mathbb{Z}\to\mathbb{Z}$, $f(n)=n+2$. Ela é injetiva porque $n_1+2=n_2+2$ implica $n_1=n_2$. É sobrejetiva porque, dado $y\in\mathbb{Z}$, o inteiro $n=y-2$ satisfaz $f(n)=y$. Logo, f é bijetiva e $f^{-1}(y)=y-2$.



\section{Problema 3 - indução}



Prove que $1+2+\cdots+n=n(n+1)/2$. A base $n=1$ é imediata. Supondo a fórmula para k,



\[1+\cdots+k+(k+1)=\frac{k(k+1)}{2}+(k+1)=\frac{(k+1)(k+2)}{2},\]



que é a fórmula para $k+1$.



\section{Problema 4 - contagem}



Quantas palavras de cinco letras distintas podem ser formadas com um alfabeto de 26 letras, contendo A mas sem A na primeira posição? Escolha a posição de A entre as quatro posições permitidas e preencha as demais com quatro letras distintas escolhidas e ordenadas entre as 25 restantes:



\[4\cdot A(25,4)=4\cdot25\cdot24\cdot23\cdot22.\]



\section{Checklist para provas e avaliações}



\begin{itemize}
\item Declare o domínio das variáveis.
\item Separe hipótese e conclusão.
\item Cite a definição usada em cada argumento central.
\item Em indução, identifique base, hipótese e passo.
\item Em relações, teste cada propriedade separadamente.
\item Em funções, não confunda contradomínio com imagem.
\item Em contagem, decida primeiro se a ordem e a repetição importam.
\item Revise se o resultado responde exatamente ao que foi perguntado.
\end{itemize}



\chapter{Apêndice - fórmulas essenciais}



\section{Lógica e conjuntos}



\[p\to q\equiv\neg p\lor q, \qquad p\to q\equiv\neg q\to\neg p.\]



\[\neg(\forall x\,P(x))\equiv\exists x\,\neg P(x), \qquad \neg(\exists x\,P(x))\equiv\forall x\,\neg P(x).\]



\[|A\cup B|=|A|+|B|-|A\cap B|.\]



\section{Contagem}



\[P_n=n!, \qquad A(n,k)=\frac{n!}{(n-k)!}, \qquad \binom{n}{k}=\frac{n!}{k!(n-k)!}.\]



\section{Fontes internas utilizadas}



Este livro foi escrito como material autoral de revisão a partir das Listas 1 a 10 de Matemática Discreta de 2025/1 e das anotações digitalizadas presentes em `resumos/originais/`. As transcrições automáticas foram usadas apenas como apoio de localização; trechos incertos não foram reproduzidos como fatos.
\end{document}
